\clearpage\thispagestyle{empty}
\pdfbookmark[0]{Quellen zum Kometenbericht (März und April 1818)}{comet1818zaa-sources}
\begin{center}\textbf{Quellen zum Kometenbericht (März und April 1818)}\end{center}
\noindent Diese KI-Übertragung ins heutige Deutsch folgt dem Bericht von Carl Friedrich Gauß in den \href{https://gdz.sub.uni-goettingen.de/id/PPN236007467}{\textit{Werken}, Band~VI} (Göttingen, 1874), S.~402--403. Der Erstabdruck erschien als Beitrag XXXIV, \textit{Auszug aus einem Briefe des Herrn Hofrath Ritter Gauß}, in der von Bernhard von Lindenau und Johann Gottlieb Friedrich Bohnenberger herausgegebenen \textit{Zeitschrift für Astronomie und verwandte Wissenschaften}, Band~V, März und April 1818, S.~\href{https://api.digitale-sammlungen.de/iiif/image/v2/bsb10061007_00286/full/1600,/0/default.jpg}{276}--\href{https://api.digitale-sammlungen.de/iiif/image/v2/bsb10061007_00287/full/1600,/0/default.jpg}{277}. Die Links führen zum Exemplar der Bayerischen Staatsbibliothek.
\par\smallskip\noindent Der Erstabdruck gibt den Logarithmus als 0,07937 und die Deklination des Vergleichssterns als $45^\circ\,9'$ an. In der Werkausgabe stehen 0.07937 und $15^\circ\,9'$. Sie druckt für den Zeitabstand des Kometen $14''98$, wo der Erstabdruck $14''{,}98$ hat. Die Werkausgabe enthält außerdem die vier späteren Beobachtungen des Sterns mit dem Vermerk „Handschrift“. KI-Übertragung und KI-Satz.
\clearpage
